\documentclass[12pt,a4paper]{article}

\usepackage[utf8]{inputenc}
\usepackage[spanish]{babel}
\usepackage[margin=2.5cm]{geometry}
\usepackage{amsmath, amssymb, amsfonts}
\usepackage{graphicx}
\usepackage{hyperref}
\usepackage{enumitem}
\usepackage{titlesec}
\usepackage{cite}
\usepackage{booktabs}
\usepackage{float}

% --- Configuración de Estilo ---
\titleformat{\section}{\large\bfseries}{\thesection.}{1em}{}
\hypersetup{colorlinks=true, linkcolor=blue, citecolor=blue, urlcolor=blue}

\title{\textbf{Proyecto de Investigación: Modelos Probabilísticos en Esports}\\ 
\Large Primera Entrega: Aproximación Bayesiana para la Predicción en CS:GO}
\author{Julian Jimenez y equipo}
\date{\today}

\begin{document}

\maketitle

\section{Introducción y Justificación}
El presente proyecto aborda la estimación de la probabilidad de victoria en encuentros profesionales de \textit{Counter-Strike: Global Offensive} (CS:GO) utilizando inferencia bayesiana. A diferencia de los modelos frecuentistas tradicionales que se centran en la clasificación binaria (Accuracy), nuestro enfoque prioriza la \textbf{calibración probabilística} y la \textbf{cuantificación de la incertidumbre epistémica}.

\section{Síntesis de la Literatura}
Basado en la revisión sistemática (Marco PICOC), se han identificado patrones críticos en la analítica de esports:
\begin{enumerate}
    \item \textbf{Sesgo de Estacionariedad:} Los modelos actuales suelen ignorar la volatilidad de los \textit{rosters}.
    \item \textbf{Falta de Calibración:} Los modelos de Deep Learning a menudo sufren de exceso de confianza (\textit{overconfidence}).
    \item \textbf{Uso de Métricas Inadecuadas:} Predomina el Accuracy sobre el Brier Score o Log-Loss.
\end{enumerate}

\section{Metodología Propuesta: Modelo Bayesiano Jerárquico}
Se propone el uso de una \textbf{Regresión Logística Bayesiana Jerárquica} implementada en \texttt{PyMC}.

\subsection{Definición de Priors}
Utilizaremos distribuciones de probabilidad para los coeficientes de:
\begin{itemize}
    \item Diferencial de ranking (HLTV).
    \item Rendimiento histórico por mapa (\textit{Map Pool}).
    \item Factor de forma reciente (\textit{Momentum}).
\end{itemize}

Se explorará el uso de \textit{opening odds} (cuotas de apertura) como \textbf{priors informativas} para mejorar la convergencia del modelo MCMC.

\section{Primera Aproximación: Análisis Exploratorio de Datos (EDA)}
En esta fase inicial, se ha implementado el sistema de captura de datos automático que consolida información de HLTV y Kaggle. El dataset cuenta con:
\begin{itemize}
    \item Resultados históricos de más de 10,000 partidas.
    \item Datos detallados de economía y vetos de mapas.
\end{itemize}

\section{Cronograma y Trabajo Futuro}
Las próximas semanas se centrarán en:
\begin{enumerate}
    \item Limpieza y transformación de datos (\textit{Feature Engineering}).
    \item Definición formal de hiperparámetros en el modelo jerárquico.
    \item Pruebas iniciales de muestreo utilizando el algoritmo NUTS.
\end{enumerate}

\bibliographystyle{plain}
\begin{thebibliography}{99}
    \bibitem{wang2023} Wang, L., \& Davis, K. (2023). \textit{Bayesian Logistic Regression for Pre-match Prediction}.
    \bibitem{smith2020} Smith, J., \& Brown, A. (2020). \textit{Predictive analytics in esports}.
\end{thebibliography}

\end{document}
